% TOP_PROBLEM: {"id": 30006986, "title": "Congruence Subgroup Property for Mapping Class Groups of Surfaces of Genus at Least Three", "category_id": 17, "category": {"id": 17, "name": "group_theory", "display_name": "Group Theory", "description": "Problems about groups, group actions, representations, and related algebraic structures.", "slug": "group-theory", "order_index": 17, "created_at": "2026-07-31T15:26:25.671Z"}, "status": "open"}
% FIRST_POSTED: 2026-09-27
% !TeX program = lualatex
% ATTEMPT_STATUS: unresolved
\documentclass[11pt]{article}
\usepackage[margin=25mm]{geometry}
\usepackage{fontspec}
\setmainfont{Latin Modern Roman}
\usepackage{amsmath,amssymb,mathtools}
\usepackage{unicode-math}
\setmathfont{Latin Modern Math}
\usepackage{xurl,hyperref}
\hypersetup{colorlinks=true,urlcolor=blue}
\setcounter{secnumdepth}{0}
\setlength{\parindent}{0pt}
\setlength{\parskip}{5pt}
\setlength{\emergencystretch}{3em}
\begin{document}
\section*{\texorpdfstring{Congruence Subgroup Property for Mapping Class Groups of Surfaces of Genus at Least Three}{Congruence Subgroup Property for Mapping Class Groups of Surfaces of Genus at Least Three}}\label{30006986}
\subsection{Definitions and mathematical statement}
For a closed connected orientable surface $S$ of genus $g\ge3$, let $\operatorname{Mod}(S)$ be the group of orientation-preserving diffeomorphisms modulo isotopy. A characteristic subgroup $N\triangleleft\pi_1(S)$ is preserved by every automorphism, so the outer action of the mapping class group descends to $\pi_1(S)/N$. For every finite-index subgroup $H\le\operatorname{Mod}(S)$, the conjecture asks for a finite-index characteristic $N$ such that
\[
 \ker\!\left(\operatorname{Mod}(S)\longrightarrow
 \operatorname{Out}(\pi_1(S)/N)\right)\subseteq H.
\]
The finite quotient may depend on $H$. It is not enough merely to separate individual mapping classes by some finite quotient: the displayed containment must hold for every finite-index subgroup.
\par\smallskip\textit{Formulation and concept references:} \href{https://arxiv.org/pdf/math/0608325}{[S1]}.

\subsection{Short English statement}
Does every finite-index subgroup of a higher-genus mapping class group contain the kernel of its action on a suitable finite characteristic quotient of the surface group?
\subsection{Research attempt: cofinality of characteristic congruence kernels}
\textbf{Status and scope.} This attempt does not prove the congruence subgroup property. The quantifier over all finite-index subgroups is essential: separating individual mapping classes by finite quotients is insufficient. The source check on 27 September 2026 found Boggi's primary manuscript [E1], which proves a congruence property for specified symmetric mapping class groups, under a genus condition on the quotient surface. That result is not asserted here for the full mapping class group of every closed surface of genus at least three. The arguments below isolate exactly what an elementary finite-quotient approach would still need.

Write $\Pi=\pi_1(S)$ and $\Gamma=\operatorname{Mod}(S)$. For a finite-index characteristic subgroup $N$ of $\Pi$, write
\[
 C_N=\ker\bigl(\Gamma\longrightarrow\operatorname{Out}(\Pi/N)\bigr).
\]
This use of outer automorphisms matters. A mapping class gives an automorphism of the fundamental group only after a basepoint choice, and changing that choice introduces an inner automorphism. Requiring the induced automorphism to be literally the identity would define a different condition without an additional marking.

\textbf{Lemma 343.1 (characteristic refinements exist).} If a group $P$ is generated by $d$ elements, every subgroup $U$ of finite index $m$ contains a finite-index characteristic subgroup of $P$.

\emph{Proof.} There are finitely many subgroups of index at most $m$. Indeed, an index-$j$ subgroup occurs as a point stabilizer in a permutation action on $j$ letters, and a homomorphism to the symmetric group is determined by the images of the $d$ generators. The bound $j(j!)^d$ on the number of possibilities is more than sufficient. Intersect all subgroups of index at most $m$, obtaining $N$. Automorphisms permute this finite collection, so $N$ is characteristic. A finite intersection of finite-index subgroups has finite index: the map from $P/N$ into the product of the corresponding coset sets is injective as a map of sets. Finally, $N\subseteq U$. $\square$

Thus there is no scarcity of finite characteristic quotients of $\Pi$. The unresolved point is whether their actions detect every finite quotient of $\Gamma$ in the stronger, uniform sense required by the question.

If $N\subseteq N'$ are characteristic, then
\[
 C_N\subseteq C_{N'}.
\]
An automorphism which is inner modulo $N$ remains inner after passage to the quotient modulo $N'$. Moreover, $N_1\cap N_2$ is again characteristic of finite index, and
\[
 C_{N_1\cap N_2}\subseteq C_{N_1}\cap C_{N_2}.
\]
Equality in this last display is not needed and is not claimed. These containments show that the congruence kernels form a basis of normal finite-index subgroups for a group topology on $\Gamma$.

\textbf{Lemma 343.2 (the exact finite-quotient reduction).} The selected assertion is equivalent to the following statement: for every homomorphism $q:\Gamma\to F$ to a finite group, there is a finite-index characteristic $N\lhd\Pi$ such that $q$ factors through the image of $\Gamma$ in $\operatorname{Out}(\Pi/N)$.

\emph{Proof.} The factorization exists exactly when $C_N\subseteq\ker q$. This proves one direction by applying the selected assertion to $\ker q$. Conversely, if $H\leq\Gamma$ has index $m$, the action on $\Gamma/H$ has finite image and kernel
\[
 L=\bigcap_{\gamma\in\Gamma}\gamma H\gamma^{-1}\subseteq H.
\]
The intersection has at most finitely many distinct conjugate factors and $[\Gamma:L]\leq m!$. Apply the finite-quotient statement to $\Gamma\to\Gamma/L$. Then $C_N\subseteq L\subseteq H$. $\square$

This reduction makes clear that the problem concerns arbitrary finite actions of the mapping class group, including ones not initially presented through the topology of a surface cover.

\textbf{A completion formulation and a failed shortcut.} Let $\widehat\Gamma$ be the ordinary profinite completion, and let
\[
 \overline\Gamma^{\,c}=\varprojlim_N\Gamma/C_N
\]
be the completion defined by these congruence kernels. There is a natural continuous surjection $\widehat\Gamma\to\overline\Gamma^{\,c}$. To see surjectivity, its image is compact and hence closed, while the diagonal image of $\Gamma$ is dense in the inverse limit and lies in the image. Density follows by refining any finite collection of the kernels by one kernel of the displayed family.

The selected assertion is equivalent to injectivity of this surjection. If the congruence kernels are cofinal among all normal finite-index subgroups, the two inverse-limit systems give the same completion. Conversely, a continuous bijection between these compact Hausdorff groups is a homeomorphism. The open kernel associated with any finite quotient of $\Gamma$ then contains the inverse image of a basic congruence neighborhood; restriction to $\Gamma$ gives $C_N\subseteq\ker q$.

This is stronger than asking whether $\bigcap_N C_N=1$ inside $\Gamma$. A completely explicit comparison is $\Gamma=\mathbb Z$ with the family $p^k\mathbb Z$, for a fixed prime $p$. Their intersection is zero, so they separate every pair of distinct integers. Nevertheless, for a different prime $q$ no $p^k\mathbb Z$ lies inside $q\mathbb Z$. The finite quotient modulo $q$ never factors through a quotient modulo a power of $p$. The map $\widehat{\mathbb Z}\to\mathbb Z_p$ has a nontrivial kernel even though the original integers inject into $\mathbb Z_p$. A proof that only separates nonidentity mapping classes could therefore miss the entire obstruction.

\textbf{An explicit family for which the required inclusion holds.} For an integer $m\geq2$, let
\[
 N_m=[\Pi,\Pi]\Pi^m.
\]
Here $\Pi^m$ denotes the subgroup generated by all $m$th powers. Both factors are characteristic, and
\[
 \Pi/N_m\cong H_1(S;\mathbb Z/m\mathbb Z)
       \cong(\mathbb Z/m\mathbb Z)^{2g}.
\]
This quotient is finite and abelian, so its outer automorphism group is its automorphism group. Consequently $C_{N_m}$ is exactly the kernel of the action on first homology modulo $m$. Every finite-index $H$ which contains this homological level subgroup satisfies the desired conclusion, with the explicitly specified $N_m$.

The calculation also gives a useful test on powers of Dehn twists. Choose a symplectic basis for integral homology and let $v$ be the primitive homology class of a nonseparating simple closed curve. With one fixed orientation convention, its twist acts by
\[
 T_v(w)=w+\langle w,v\rangle v,
 \qquad T_v^k(w)=w+k\langle w,v\rangle v.
 \tag{343.1}
\]
The second formula follows by induction, since $\langle v,v\rangle=0$. As $v$ is primitive, there is an integral $w$ with $\langle w,v\rangle=1$. It follows that
\[
 T_v^k\equiv I\pmod m\quad\Longleftrightarrow\quad m\mid k.
\]
For the reverse implication in this equivalence, applying the operator to that $w$ gives $kv=0$ modulo $m$; the coordinates of a primitive integral vector generate the unit ideal, so $m\mid k$. This argument works for composite $m$ as well as prime $m$.

Equation (343.1) also exposes the limitation of using only abelian characteristic quotients. A separating twist has homology class zero for its twisting curve, and its homological action is trivial. Thus this entire family of quotients ignores the Torelli subgroup. The full problem allows nonabelian finite characteristic quotients of $\Pi$ precisely because the homological information alone need not control a general finite-index $H$.

\textbf{Why enumeration does not finish the proof.} Finite generation of $\Pi$ allows characteristic refinements to be built from finite permutation actions. One can also enumerate finite presentations for finite quotients and test many compatible actions in examples. None of this supplies a bound, in terms of $[\Gamma:H]$, on the complexity of a characteristic quotient whose kernel lies in $H$. More fundamentally, the existence of such a quotient is the conclusion still missing. An exhaustive search would be a semidecision procedure only after the relevant group-action computations were made effective; termination would still require the conjecture.

The exact finite diagnostics accompanying this attempt verify the symplectic transvection identity and its order modulo several small prime and composite moduli. They also check the divisibility obstruction for the $p$-power filtration of $\mathbb Z$. These are checks of the reductions and examples, not tests of cofinality for the mapping class group.

\textbf{Outcome and first gap.} Characteristic finite quotients form a directed family, arbitrary finite-index subgroups reduce to finite normal quotients, and the desired conclusion holds for every subgroup containing an explicit homological level. The missing assertion is that every finite quotient of $\operatorname{Mod}(S)$ is controlled by some member of the much larger nonabelian characteristic family. Neither element separation, homological transvections, nor the special symmetric-group theorem in [E1] proves that assertion. The attempt remains unresolved.

\subsection{Sources}
\begin{itemize}
\item[S1]Nikolai V. Ivanov, Fifteen problems about the mapping class groups (2006).
\url{https://arxiv.org/pdf/math/0608325}.
\textit{Formulation source.}
\item[E1]Marco Boggi, \emph{A congruence subgroup property for symmetric mapping class groups}, arXiv:2408.12486 (2024).
\url{https://arxiv.org/abs/2408.12486}.
\textit{Primary source consulted 27 September 2026; the quotient-surface genus restriction is retained in the scope discussion.}
\end{itemize}
\end{document}
